\chapter{DSL Parsing}
\label{ch:dsl}

The DSL is a small arrow-chain grammar designed to be read aloud
(``User authenticates to a Host and executes a Process'') and to
fit on one line of a UI input. The parser is a hand-written
single-pass recursive-descent walker --- no grammar file, no
parser-generator dependency.

\section{Grammar}

\begin{verbatim}
hypothesis      := step ( step )* modifier*
step            := entity_type arrow
arrow           := "-[" rel_type predicates? "]->" entity_type binding?
                   predicates? dest_predicates?
entity_type     := ident ( "(" binding ")" )? | "*"
rel_type        := ident | "*"
predicates      := "{" predicate ( "," predicate )* "}"
predicate       := key ("=" | "~" | " in ") (literal | list)
modifier        := having | within | k_clause
having          := "HAVING" agg_op "(" agg_arg ")" cmp literal
within          := "WITHIN" int time_unit
k_clause        := "{k=" int "}"
\end{verbatim}

Every token type is recognised by a character-class probe, so the
parser performs exactly one pass over the input string. There is no
backtracking; the grammar is LL(1) after lexer-level lookahead for
predicates.

\section{Algorithm}

\begin{algorithm}
\caption{\textsc{Parse-Dsl}}\label{alg:parse-dsl}
\KwIn{string $\texttt{src}$, optional hypothesis name $\texttt{nm}$}
\KwOut{\code{Hypothesis} or \code{DslError}}
$\text{toks} \gets $ \textsc{Tokenize}($\texttt{src}$) \Comment*[r]{\Ord{|\texttt{src}|}}
$\text{steps} \gets \emptyset$\;
$\sigma_0 \gets $ \textsc{Parse-Entity-Type}(toks)\;
\While{toks has remaining ``$\to$''}{
    $\rho, P^{\text{edge}} \gets $ \textsc{Parse-Edge-Spec}(toks)\;
    $\sigma_i \gets $ \textsc{Parse-Entity-Type}(toks)\;
    $P^{\text{dst}} \gets $ \textsc{Parse-Dest-Predicates}(toks)\;
    push \code{HypothesisStep}$(\sigma_{i-1}, \rho, \sigma_i, P^{\text{edge}}, P^{\text{dst}})$ onto steps\;
}
$k \gets$ parse optional $\texttt{\{k=N\}}$, default 1\;
$A \gets$ \textsc{Parse-Having}(toks)\;
$w \gets$ \textsc{Parse-Within}(toks)\;
\Return \code{Hypothesis}(nm, steps, $k$, $A$, $w$)\;
\end{algorithm}

\begin{complexity}
$\Ord{|\texttt{src}|}$ time and space. Every token is
visited exactly once; the \code{HashMap}-based predicate lists are
$\Ord{1}$ to populate. Space is linear in the parsed chain length
plus the predicate literals.
\end{complexity}

\begin{theorem}[Parser determinism]
\textsc{Parse-Dsl} is deterministic: for any fixed input string, it
produces the same \code{Hypothesis} or the same \code{DslError} on
every invocation.
\end{theorem}

\begin{proof}
No hash-map iteration order, no random values, no calls to the
system clock anywhere in the parser. All data structures built
during parsing (the final \code{edge\_predicates} and
\code{dest\_predicates} on each step) preserve insertion order.
\qed
\end{proof}

\begin{inthecode}
\filepath{graph\_hunter\_core/src/dsl.rs:47} ---
\code{parse\_dsl} entry point.\\
\filepath{graph\_hunter\_core/src/dsl.rs} --- full grammar
implementation (private helpers).
\end{inthecode}
